\chapter{Teoría de la medida}\label{ch:b3:measure}

¿Cuánto mide un subconjunto de $\R$? La respuesta ingenua --- asignar a
todo conjunto una longitud invariante por traslaciones que extienda la
de los intervalos --- es \emph{imposible}: la construcción de Vitali, al
final de este capítulo, produce un conjunto sin longitud coherente. La
teoría de la \omterm{def:b3:measure:measure}{medida} es la retirada
disciplinada: restringir la atención a una clase amplia de conjuntos
\emph{medibles}, sobre la que existe una longitud numerablemente
aditiva, y es única. Las recompensas son inmensas --- la integral de
Lebesgue (el~\cref{ch:b3:lebesgue}), los espacios $L^p$ del análisis
funcional y toda la probabilidad moderna (el~\cref{ch:b3:probability})
se apoyan en los tres teoremas que aquí se demuestran: el lema de
unicidad de Dynkin, el teorema de extensión de Carathéodory y la
existencia de la \omterm{def:b3:measure:lebesgueouter}{medida de
Lebesgue}.

\section{$\sigma$-álgebras}

\begin{definition}\label{def:b3:measure:sigmaalgebra}
Una \emph{$\sigma$-álgebra}\index{sigma-algebra@$\sigma$-álgebra} sobre
un conjunto $X$ es una familia $\mathcal A$ de subconjuntos que contiene
$\varnothing$ y es estable por complementación y por uniones
\emph{numerables} (y, por tanto, por intersecciones numerables y
diferencias de conjuntos, y contiene $X$). El par $(X, \mathcal A)$ es
un \emph{espacio medible}; los elementos de $\mathcal A$ son los
\emph{conjuntos medibles}. Para una familia $\mathcal E$ de
subconjuntos, $\sigma(\mathcal E)$ denota la menor $\sigma$-álgebra que
contiene $\mathcal E$ (la intersección de todas ellas --- una
intersección de $\sigma$-álgebras lo es).
\end{definition}

\begin{definition}\label{def:b3:measure:borel}
La \emph{$\sigma$-álgebra de Borel}\index{sigma-algebra de
Borel@$\sigma$-álgebra de Borel} de un espacio
\omterm{def:b3:topology:topology}{topológico} es
$\mathcal B(X) = \sigma(\{\text{abiertos}\})$. Sobre $\R$:
$\mathcal B(\R)$ está también generada por los intervalos abiertos, por
los intervalos cerrados, por las semirrectas $(-\infty, a]$ y por las
semirrectas con extremo racional (el~\cref{exo:b3:measure:1}) --- cada
familia genera los abiertos mediante operaciones numerables; por
ejemplo, todo abierto de $\R$ es
unión numerable de intervalos abiertos con datos racionales.
\end{definition}

\begin{definition}\label{def:b3:measure:dynkin}
Un \emph{$\pi$-sistema}\index{pi-sistema@$\pi$-sistema} es una familia
estable por intersecciones finitas. Un
\emph{$\lambda$-sistema}\index{lambda-sistema@$\lambda$-sistema} (clase
de Dynkin) es una familia $\mathcal D$ tal que: $X \in \mathcal
D$; $A, B \in \mathcal D$, $A \subseteq B$ $\Rightarrow$ $B
\setminus A \in \mathcal D$; y $A_n \uparrow A$, $A_n \in
\mathcal D$ $\Rightarrow$ $A \in \mathcal D$.
\end{definition}

\begin{theorem}[Lema $\pi$--$\lambda$ de Dynkin]
\label{thm:b3:measure:dynkin}
Si un $\lambda$-sistema $\mathcal D$ contiene un $\pi$-sistema
$\mathcal P$, entonces $\mathcal D \supseteq
\sigma(\mathcal P)$.\index{lema de Dynkin}
\end{theorem}

\begin{proof}
Sea $\mathcal D_0$ el menor $\lambda$-sistema que contiene $\mathcal P$
(intersección de todos ellos); basta ver que $\mathcal D_0$ es una
$\sigma$-álgebra, pues entonces
$\sigma(\mathcal P) \subseteq \mathcal D_0 \subseteq \mathcal
D$. Un $\lambda$-sistema estable por intersecciones finitas \emph{es}
una $\sigma$-álgebra: complementos ($X \setminus A = X
\setminus A$ con $A \subseteq X$), uniones finitas ($A \cup B =
X \setminus ((X\setminus A)\cap(X\setminus B))$) y uniones numerables
mediante $\bigcup_{k \leq n}A_k \uparrow \bigcup_kA_k$. Demostremos,
pues, que $\mathcal D_0$ es un $\pi$-sistema, en dos pasos. Sea
\[
\mathcal D_1 = \{A \in \mathcal D_0 : A \cap P \in \mathcal
D_0 \ \forall P \in \mathcal P\}.
\]
$\mathcal D_1$ es un $\lambda$-sistema (los tres axiomas se verifican
intersecando con $P$: por ejemplo, $(B\setminus A)\cap P
= (B \cap P)\setminus(A \cap P)$, una diferencia propia dentro de
$\mathcal D_0$) y contiene $\mathcal P$ ($\pi$-sistema):
$\mathcal D_1 = \mathcal D_0$. Sea ahora
\[
\mathcal D_2 = \{A \in \mathcal D_0 : A \cap D \in \mathcal
D_0\ \forall D \in \mathcal D_0\}.
\]
Por el paso anterior, $\mathcal D_2 \supseteq \mathcal P$; y
$\mathcal D_2$ es un $\lambda$-sistema por la misma verificación:
$\mathcal D_2 = \mathcal D_0$, lo que dice exactamente que
$\mathcal D_0$ es estable por intersecciones.
\end{proof}

\section{Medidas}

\begin{definition}\label{def:b3:measure:measure}
Una \emph{medida}\index{medida} sobre $(X, \mathcal A)$ es una
aplicación $\mu \colon \mathcal A \to [0, +\infty]$ con
$\mu(\varnothing)
= 0$ que es \emph{$\sigma$-aditiva}: para conjuntos disjuntos dos a dos
$(A_n)_{n\in\N}$,
\[
\mu\Bigl(\bigsqcup_n A_n\Bigr) = \sum_n \mu(A_n).
\]
$(X, \mathcal A, \mu)$ es un \emph{espacio de medida}; $\mu$ es
\emph{finita} si $\mu(X) < \infty$, una \emph{medida de probabilidad} si
$\mu(X) = 1$, y \emph{$\sigma$-finita} si $X$ es unión numerable de
conjuntos de medida finita. Ejemplos: la medida de conteo sobre
$(\N, \mathcal P(\N))$; la masa de Dirac $\delta_a(A)
= \mathbf 1_{a \in A}$; y, objeto de este capítulo, la
\omterm{def:b3:measure:lebesgueouter}{medida de Lebesgue}.
\end{definition}

\begin{proposition}\label{prop:b3:measure:basics}
Sea $\mu$ una \omterm{def:b3:measure:measure}{medida}. (a) Monotonía:
$A \subseteq B
\Rightarrow \mu(A) \leq \mu(B)$. (b) Subaditividad numerable:
$\mu(\bigcup A_n) \leq \sum\mu(A_n)$. (c)
\omterm{def:b3:topology:continuity}{Continuidad} por abajo:
$A_n \uparrow A \Rightarrow \mu(A_n) \to \mu(A)$. (d)
\omterm{def:b3:topology:continuity}{Continuidad} por arriba:
$A_n \downarrow A$ \emph{con
$\mu(A_1) < \infty$} $\Rightarrow \mu(A_n) \to \mu(A)$.
\end{proposition}

\begin{proof}
(a) $B = A \sqcup (B\setminus A)$. (b) Disjúntese: los $B_n = A_n
\setminus \bigcup_{k<n}A_k$ son disjuntos con la misma unión, y
$\mu(B_n) \leq \mu(A_n)$. (c) $A = \bigsqcup_n (A_n
\setminus A_{n-1})$ ($A_0 = \varnothing$): las sumas parciales de
$\sum\mu(A_n\setminus A_{n-1})$ son $\mu(A_n)$. (d) Aplíquese (c) a
$A_1 \setminus A_n \uparrow A_1 \setminus A$ y réstese de $\mu(A_1)$ ---
la finitud legitima la resta. Contraejemplo sin ella:
$A_n = [n, \infty)$ para la \omterm{def:b3:measure:lebesgueouter}{medida
de Lebesgue}: $A_n \downarrow \varnothing$ pero $\mu(A_n) =
\infty$.
\end{proof}

\begin{theorem}[Unicidad]\label{thm:b3:measure:uniqueness}
Sean $\mu, \nu$ \omterm{def:b3:measure:measure}{medidas} sobre
$\sigma(\mathcal P)$ y sea $\mathcal
P$ un $\pi$-sistema con $\mu = \nu$ sobre $\mathcal P$. Si existen
conjuntos $P_k \in \mathcal P$ con $P_k \uparrow X$ y
$\mu(P_k) < \infty$, entonces $\mu = \nu$ en todo
$\sigma(\mathcal P)$.
\end{theorem}

\begin{proof}
Fíjese $k$ y considérense las \omterm{def:b3:measure:measure}{medidas}
finitas $\mu_k(A) = \mu(A \cap
P_k)$ y $\nu_k(A) = \nu(A \cap P_k)$ sobre $\sigma(\mathcal
P)$: coinciden sobre $\mathcal P$, pues $P \cap P_k \in \mathcal
P$ ($\pi$-sistema), y asignan a $X$ el mismo valor finito $\mu(P_k)$. La
clase $\mathcal D = \{A : \mu_k(A) =
\nu_k(A)\}$ es un $\lambda$-sistema: $X \in \mathcal D$; las diferencias
propias, por resta (valores finitos); los límites crecientes, por
\omterm{def:b3:topology:continuity}{continuidad} por abajo
(la~\cref{prop:b3:measure:basics}(c)). Contiene el $\pi$-sistema
$\mathcal P$, de modo que Dynkin (el~\cref{thm:b3:measure:dynkin}) da
$\mathcal D \supseteq
\sigma(\mathcal P)$: $\mu_k = \nu_k$ en todas partes. Por último, para
$A \in \sigma(\mathcal P)$ cualquiera, la
\omterm{def:b3:topology:continuity}{continuidad} por abajo a lo largo de
$A
\cap P_k \uparrow A$ da $\mu(A) = \lim_k\mu_k(A) =
\lim_k\nu_k(A) = \nu(A)$.
\end{proof}

\section{Medidas exteriores y el teorema de Carathéodory}

\begin{definition}\label{def:b3:measure:outer}
Una \emph{medida exterior}\index{medida exterior} sobre $X$ es una
aplicación $\mu^* \colon \mathcal P(X) \to [0, \infty]$ con
$\mu^*(\varnothing) = 0$, monótona y numerablemente subaditiva. Un
conjunto $A$ es \emph{$\mu^*$-medible} (Carathéodory) si divide
aditivamente todo conjunto:
\[
\mu^*(E) = \mu^*(E \cap A) + \mu^*(E \setminus A)
\qquad \text{para todo } E \subseteq X
\]
($\leq$ se cumple siempre por subaditividad; el contenido es $\geq$).
\end{definition}

\begin{theorem}[Carathéodory]\label{thm:b3:measure:caratheodory}
Los conjuntos $\mu^*$-medibles forman una $\sigma$-álgebra $\mathcal
M$, y $\mu^*\restriction_{\mathcal M}$ es una
\omterm{def:b3:measure:measure}{medida}. Además, todo conjunto con
$\mu^*(N) = 0$ pertenece a $\mathcal M$ (la
\omterm{def:b3:measure:measure}{medida} es
\emph{\omterm{def:b3:complete:complete}{completa}}).\index{teorema de
Carathéodory}
\end{theorem}

\begin{proof}
$\mathcal M$ contiene $\varnothing$ y es estable por complementación (la
condición que la define es simétrica en $A$, $X\setminus A$).
\emph{Uniones finitas}: sean $A, B \in \mathcal
M$ y $E$ arbitrario; dividiendo $E$ por $A$ y después cada trozo por
$B$:
\[
\mu^*(E) = \mu^*(E\cap A\cap B) + \mu^*(E\cap A\setminus B) +
\mu^*(E\cap B\setminus A) + \mu^*(E\setminus(A\cup B)).
\]
Los tres primeros trozos recubren $E \cap (A \cup B)$, de modo que la
subaditividad da $\mu^*(E) \geq \mu^*(E\cap(A\cup B)) +
\mu^*(E\setminus(A\cup B))$: $A \cup B \in \mathcal M$. Por inducción,
las uniones finitas; y con los complementos, quedan disponibles las
manipulaciones de disjunción finita.

\emph{Aditividad sobre $\mathcal M$}: para $A, B \in
\mathcal M$ disjuntos y $E$ cualquiera:
$\mu^*(E\cap(A\sqcup B)) = \mu^*(E\cap
A) + \mu^*(E \cap B)$ (divídase por $A$); por inducción,
\[
\mu^*\Bigl(E \cap \bigsqcup_{k\leq n}A_k\Bigr)
= \sum_{k\leq n}\mu^*(E\cap A_k).
\tag{$*$}
\]

\emph{Uniones numerables}: sean $(A_k) \subseteq \mathcal M$ disjuntos
(basta con ello, por disjunción dentro del álgebra $\mathcal M$),
$A = \bigsqcup A_k$, $E$ arbitrario. Usando
$\bigsqcup_{k\leq n}A_k \in \mathcal M$ y la monotonía:
\[
\mu^*(E) = \mu^*\Bigl(E\cap\bigsqcup_{k\leq n}A_k\Bigr) +
\mu^*\Bigl(E\setminus\bigsqcup_{k\leq n}A_k\Bigr)
\geq \sum_{k \leq n}\mu^*(E\cap A_k) + \mu^*(E\setminus A)
\]
por ($*$). Hágase $n \to \infty$ y úsese la subaditividad numerable en
sentido contrario:
\[
\mu^*(E) \geq \sum_{k}\mu^*(E\cap A_k) + \mu^*(E\setminus A)
\geq \mu^*(E \cap A) + \mu^*(E\setminus A) \geq \mu^*(E):
\]
todas las desigualdades son igualdades. Esto demuestra a la vez $A \in
\mathcal M$ y, tomando $E = A$, la aditividad numerable de $\mu^*$ sobre
$\mathcal M$.

\emph{Conjuntos nulos}: si $\mu^*(N) = 0$, entonces, para cualquier $E$:
$\mu^*(E \cap N) + \mu^*(E\setminus N) \leq 0 + \mu^*(E)$:
$N \in \mathcal M$.
\end{proof}

\section{La medida de Lebesgue en \texorpdfstring{$\R$}{R}}

\begin{definition}\label{def:b3:measure:lebesgueouter}
La \emph{medida exterior de Lebesgue}\index{medida de Lebesgue} de $A
\subseteq \R$ es
\[
\lambda^*(A) = \inf\Bigl\{\sum_{n} (b_n - a_n) :
A \subseteq \bigcup_n \intoo{a_n}{b_n}\Bigr\}
\]
(recubrimientos numerables por intervalos abiertos).
\end{definition}

\begin{lemma}\label{lem:b3:measure:outerbasics}
$\lambda^*$ es una \omterm{def:b3:measure:outer}{medida exterior},
invariante por traslaciones, y $\lambda^*(I) = \ell(I)$ (la longitud)
para todo intervalo
$I$.
\end{lemma}

\begin{proof}
\omterm{def:b3:measure:outer}{Medida exterior}: $\varnothing$ se recubre
por intervalos arbitrariamente pequeños; la monotonía es clara;
subaditividad: dados recubrimientos de cada $A_n$ hasta
$\varepsilon 2^{-n}$ del ínfimo, su unión recubre $\bigcup A_n$ con
longitud total $\leq \sum
\lambda^*(A_n) + \varepsilon$. Invariancia por traslaciones: tradúzcanse
los recubrimientos.

Longitud: basta tratar $I = \intcc ab$ (los demás tipos difieren en los
extremos, de \omterm{def:b3:measure:outer}{medida exterior} $0$:
recúbranse por intervalos diminutos; después, compárense por encaje los
tipos $\intcc{a+\varepsilon}{b -
\varepsilon} \subseteq \intoo ab$). $\lambda^*(\intcc ab) \leq b - a$:
recúbrase por $\intoo{a-\varepsilon}{b+\varepsilon}$. Recíprocamente,
sea $\intcc
ab \subseteq \bigcup_n\intoo{a_n}{b_n}$: por \emph{\omterm{def:b3:topology:compact}{compacidad}}
(Borel--Lebesgue, el~\cref{thm:b3:topology:metriccompact}) bastan un
número finito de intervalos, digamos $I_1, \dots, I_N$. Veamos que
$\sum_{k\leq N}(b_k -
a_k) \geq b - a$ por inducción sobre $N$: elíjase $I_{k_1} \ni a$; si
$b_{k_1} > b$, hemos terminado ($b_{k_1} - a_{k_1} > b - a$); en caso
contrario, el segmento $\intcc{b_{k_1}}b$ queda recubierto por los
$N - 1$ intervalos restantes, y la inducción da
$\sum_{k \neq k_1}(b_k - a_k)
\geq b - b_{k_1}$, mientras que $b_{k_1} - a_{k_1} > b_{k_1} - a$:
súmese.
\end{proof}

\begin{theorem}[Medida de Lebesgue]\label{thm:b3:measure:lebesgue}
Todo conjunto de Borel de $\R$ es $\lambda^*$-medible. La restricción
$\lambda$ de $\lambda^*$ a la $\sigma$-álgebra
$\mathcal L = \mathcal M_{\lambda^*} \supseteq \mathcal B(\R)$
(la \emph{$\sigma$-álgebra de Lebesgue}) es la única
\omterm{def:b3:measure:measure}{medida} sobre $\mathcal B(\R)$ que
asigna a cada intervalo su longitud; es invariante por traslaciones y
$\sigma$-finita.
\end{theorem}

\begin{proof}
Por el~\cref{thm:b3:measure:caratheodory} basta ver que cada semirrecta
$A = \intoo{-\infty}c$ es $\lambda^*$-medible (las semirrectas generan
$\mathcal B$, la~\cref{def:b3:measure:borel}). Sea $E \subseteq \R$ con
$\lambda^*(E) < \infty$ y sea $\bigcup I_n \supseteq E$ un recubrimiento
con $\sum\ell(I_n) \leq \lambda^*(E) + \varepsilon$. Cada $I_n$ se
descompone en los dos intervalos $I_n' = I_n \cap A$ y
$I_n'' = I_n\setminus A$ (un intervalo menos una semirrecta es un
intervalo) con $\ell(I_n') + \ell(I_n'') = \ell(I_n)$; los $I_n'$
recubren $E \cap A$ y los $I_n''$ recubren $E \setminus A$ (agrándese
cada uno hasta un intervalo abierto de longitud $\ell +
\varepsilon2^{-n}$ para no salirse de la definición), de modo que
\[
\lambda^*(E\cap A) + \lambda^*(E\setminus A)
\leq \sum_n\bigl(\ell(I_n') + \ell(I_n'')\bigr) + 2\varepsilon
\leq \lambda^*(E) + 3\varepsilon .
\]
Unicidad: dos \omterm{def:b3:measure:measure}{medidas} que coinciden con
la longitud sobre el $\pi$-sistema de los intervalos $\intoc ab$
(finitas sobre ellos) coinciden en
$\sigma(\text{intervalos}) = \mathcal B$ por
el~\cref{thm:b3:measure:uniqueness} con $P_k = \intoc{-k}k$.
$\sigma$-finitud: $\R = \bigcup(-k, k]$.
\end{proof}

\begin{theorem}[Regularidad]\label{thm:b3:measure:regularity}
Para todo $A \in \mathcal L$:\index{regularidad de la medida de
Lebesgue}
\[
\lambda(A) = \inf\{\lambda(U) : U \supseteq A \text{ abierto}\}
= \sup\{\lambda(K) : K \subseteq A \text{ compacto}\}.
\]
\end{theorem}

\begin{proof}
\emph{Exterior}: un recubrimiento $\bigcup I_n$ con $\sum\ell(I_n) \leq
\lambda(A) + \varepsilon$ es un
abierto $U \supseteq A$ con
$\lambda(U) \leq \lambda(A) + \varepsilon$ (subaditividad); si
$\lambda(A) = \infty$, el enunciado es trivial. \emph{\omterm{def:b3:topology:interior}{Interior}}: sea
primero $A$ acotado, $A \subseteq [-M, M]$. Tómese un abierto
$U \supseteq ([-M,M]\setminus A)$ con $\lambda(U) \leq
\lambda([-M,M]\setminus A) + \varepsilon$; entonces $K = [-M,
M]\setminus U$ es \omterm{def:b3:topology:compact}{compacto},
$K \subseteq A$, y
\[
\lambda(K) = \lambda([-M,M]) - \lambda([-M,M]\cap U)
\geq \lambda([-M,M]) - \bigl(\lambda([-M,M]) - \lambda(A) +
\varepsilon\bigr) = \lambda(A) - \varepsilon .
\]
Para $A$ general: $\lambda(A) = \lim_M\lambda(A\cap[-M,M])$
(\omterm{def:b3:topology:continuity}{continuidad} por abajo) y aplíquese
el caso acotado dentro.
\end{proof}

\begin{example}\label{ex:b3:measure:cantor}
El conjunto de Cantor (el~\cref{exo:b3:topology:10}) cumple
$\lambda(C) = 0$: $C \subseteq C_n$, unión de $2^n$ intervalos de
longitud $3^{-n}$, luego $\lambda(C) \leq (2/3)^n \to 0$. Un conjunto
nulo no numerable --- el cardinal no ve la
\omterm{def:b3:measure:measure}{medida}. Recíprocamente, los
\emph{conjuntos de Cantor gordos} (el~\cref{exo:b3:measure:5}) son
densos en ninguna parte y de \omterm{def:b3:measure:measure}{medida}
positiva: la \omterm{def:b3:topology:topology}{topología} tampoco ve la
\omterm{def:b3:measure:measure}{medida}. El problema de fin de semana
lleva esta interacción a su conclusión más llamativa: hay conjuntos
medibles Lebesgue que no son de Borel.
\end{example}

\begin{theorem}[Vitali]\label{thm:b3:measure:vitali}
No existe ninguna \omterm{def:b3:measure:measure}{medida} sobre
\emph{todos} los subconjuntos de $\R$ que sea invariante por
traslaciones y asigne a cada intervalo su longitud. En particular,
$\mathcal L \neq \mathcal P(\R)$: existen conjuntos no
medibles.\index{conjunto no medible de Vitali}
\end{theorem}

\begin{proof}
Supongamos que $\mu$ fuera una tal \omterm{def:b3:measure:measure}{medida}. Sobre $\intcc01$, considérese
la equivalencia $x \sim y \iff x - y \in \Q$; por el \emph{axioma de
elección}, tómese un representante en $\intcc01$ por cada clase: un
conjunto $V$. Para $q \in \Q\cap\intcc{-1}1$, los trasladados $V + q$
son disjuntos dos a dos (dos puntos de $V$ que difirieran en un racional
serían equivalentes y, sin embargo, representantes distintos) y
\[
\intcc01 \subseteq \bigsqcup_{q \in \Q\cap\intcc{-1}1}(V + q)
\subseteq \intcc{-1}2 :
\]
la primera inclusión porque todo $x \in \intcc01$ difiere de su
representante $v$ en un racional $q = x - v \in
\intcc{-1}1$. La monotonía y la $\sigma$-aditividad dan
\[
1 \leq \sum_{q}\mu(V + q) \leq 3,
\qquad\text{con } \mu(V + q) = \mu(V) \text{ para todos } q .
\]
Una suma infinita de la constante $\mu(V)$ vale $0$ o $\infty$: ambas
cotas no pueden cumplirse a la vez. Así que no existe tal $\mu$ --- y $V
\notin \mathcal L$, pues $\lambda$ sobre $\mathcal L$ tiene todas las
propiedades usadas.
\end{proof}

\begin{remark}\label{rem:b3:measure:banachtarski}
En $\R^3$ el fracaso es más dramático: la paradoja de Banach--Tarski
descompone una bola en cinco piezas que se recomponen, mediante
rotaciones y traslaciones, en \emph{dos} bolas del mismo radio --- de
modo que ni siquiera existe un volumen finitamente aditivo e invariante
por rotaciones sobre todos los subconjuntos de $\R^3$. Las piezas son,
por supuesto, no medibles. La medibilidad no es una cautela burocrática;
es la frontera de la coherencia.
\end{remark}

\begin{method}\label{met:b3:measure:goodsets}
El \emph{principio de los buenos conjuntos}: para demostrar que todos
los conjuntos de $\sigma(\mathcal E)$ tienen una propiedad, compruébese
que los buenos conjuntos forman una $\sigma$-álgebra (o un
$\lambda$-sistema, si la propiedad es de tipo \omterm{def:b3:measure:measure}{medida} y $\mathcal E$ es
un $\pi$-sistema --- entonces, Dynkin) que contiene $\mathcal E$. Casi
todas las demostraciones de este capítulo y del siguiente son un caso
particular. Para demostrar que dos
\omterm{def:b3:measure:measure}{medidas} son iguales: compruébese sobre
un $\pi$-sistema generador más la $\sigma$-finitud
(el~\cref{thm:b3:measure:uniqueness}). Para construir una
\omterm{def:b3:measure:measure}{medida}: constrúyase una
\omterm{def:b3:measure:outer}{medida exterior} mediante recubrimientos e
invóquese a Carathéodory.
\end{method}

\section{Ejercicios}

\begin{exercise}[$\star$]\label{exo:b3:measure:1}
(a) Demostrar que $\{A \subseteq X : A$ o $X\setminus A$ es
numerable$\}$ es una $\sigma$-álgebra: la generada por los puntos. (b)
Demostrar que $\mathcal B(\R)$ está generada por cada una de estas
familias: los intervalos abiertos; los intervalos cerrados; las
semirrectas $\intoc{-\infty}a$; las semirrectas con $a \in \Q$. (c) ¿Es
la familia de las \emph{uniones finitas disjuntas de intervalos}
$\intoc ab$ una $\sigma$-álgebra? ¿Es un álgebra (estable por
complementación y uniones finitas)?
\end{exercise}

\begin{exercise}[$\star$]\label{exo:b3:measure:2}
(a) Demostrar la fórmula de inclusión--exclusión para una
\omterm{def:b3:measure:measure}{medida} finita: $\mu(A\cup
B) = \mu(A) + \mu(B) - \mu(A\cap B)$, y la versión con tres conjuntos.
(b) Dar un ejemplo que muestre que la
\omterm{def:b3:topology:continuity}{continuidad} por arriba
(la~\cref{prop:b3:measure:basics}(d)) falla sin la hipótesis de finitud.
(c) Demostrar que un conjunto numerable tiene
\omterm{def:b3:measure:lebesgueouter}{medida de Lebesgue} nula. Deducir
$\lambda(\Q) = 0$ y $\lambda(\intcc01\setminus\Q) =
1$.
\end{exercise}

\begin{exercise}[$\star\star$]\label{exo:b3:measure:3}
Sean $\mu, \nu$ \omterm{def:b3:measure:measure}{medidas} de probabilidad
sobre $\mathcal B(\R)$ con
$\mu(\intoc{-\infty}t) = \nu(\intoc{-\infty}t)$ para todo $t \in \R$.
Demostrar que $\mu = \nu$. (Esto hace de la \emph{función de
distribución} $F(t) = \mu(\intoc{-\infty}t)$ un invariante
\omterm{def:b3:complete:complete}{completo} --- el fundamento de
\cref{ch:b3:probability}.)
\end{exercise}

\begin{exercise}[$\star\star$]\label{exo:b3:measure:4}
(Borel--Cantelli, versión de \omterm{def:b3:measure:measure}{medida})
Sean $(A_n)$ medibles con $\sum_n\mu(A_n) < \infty$, y sea
$\limsup A_n =
\bigcap_N\bigcup_{n\geq N}A_n$ (los puntos que pertenecen a infinitos
$A_n$). Demostrar que $\mu(\limsup A_n) = 0$. Aplicación: para casi todo
$x \in \intcc01$, solo un número finito de $n$ cumplen
$\abs{x - p/q_n} \leq 4^{-n}$ para el $n$-ésimo racional $p/q_n$ de una
enumeración de
$\Q\cap\intcc01$.
\end{exercise}

\begin{exercise}[$\star\star$]\label{exo:b3:measure:5}
(Conjunto de Cantor gordo) Repítase la construcción de Cantor sobre
$\intcc01$, pero suprimiendo en la etapa $n$, de cada uno de los
$2^{n-1}$ intervalos, un intervalo abierto \emph{centrado} de longitud
solo $4^{-n}$. Demostrar que el $K = \bigcap K_n$ resultante es
\omterm{def:b3:topology:compact}{compacto}, tiene
\omterm{def:b3:topology:interior}{interior} vacío (no sobrevive ningún
intervalo) y
\[
\lambda(K) = 1 - \sum_{n\geq1}2^{n-1}4^{-n} = \tfrac12 :
\]
un \omterm{def:b3:topology:interior}{conjunto denso} en ninguna parte de
\omterm{def:b3:measure:measure}{medida} $\frac12$. Deducir la existencia
de un subconjunto \emph{\omterm{rem:b3:complete:meagre}{magro}} de
$\intcc01$ de \omterm{def:b3:measure:measure}{medida} total $1$, y de un
subconjunto abierto y denso de \omterm{def:b3:measure:measure}{medida}
$< \varepsilon$.
\end{exercise}

\begin{exercise}[$\star\star$]\label{exo:b3:measure:6}
Sea $\mu$ una \omterm{def:b3:measure:measure}{medida} sobre
$\mathcal B(\R)$, invariante por traslaciones, con
$c = \mu(\intoc01) < \infty$. Demostrar que $\mu =
c\,\lambda$ sobre $\mathcal B(\R)$. \emph{(Calcúlese $\mu$ sobre los
intervalos diádicos dividiendo $\intoc01$ en $2^n$ trasladados, e
invóquese después
\cref{thm:b3:measure:uniqueness}.)}
\end{exercise}

\begin{exercise}[$\star\star$]\label{exo:b3:measure:7}
(Aproximación) Sea $A \in \mathcal L$ con $\lambda(A) <
\infty$ y sea $\varepsilon > 0$. Demostrar que existe una unión
\emph{finita} de intervalos $B$ con $\lambda(A\,\triangle\,B) <
\varepsilon$ ($\triangle$ = diferencia simétrica). \emph{(Regularidad:
encájese $K \subseteq A \subseteq U$ y úsese la estructura del abierto
$U$ como unión numerable de intervalos, junto con la \omterm{def:b3:topology:compact}{compacidad} de
$K$.)}
\end{exercise}

\begin{exercise}[$\star\star\star$]\label{exo:b3:measure:8}
(Steinhaus) Sea $A \in \mathcal L$ con $\lambda(A) > 0$. Demostrar que
$A - A = \{x - y : x, y \in A\}$ contiene un intervalo alrededor de $0$.
\emph{(Redúzcase a $\lambda(A) < \infty$; por regularidad al estilo
del~\cref{exo:b3:measure:7}, hállese un intervalo $I$ con
$\lambda(A \cap I) > \frac34\ell(I)$; entonces, para $\abs t <
\frac12\ell(I)$, los conjuntos $A\cap I$ y $(A\cap I) + t$ están ambos
en un intervalo de longitud $\frac32\ell(I)$ y tienen
\omterm{def:b3:measure:measure}{medida} total $> \frac32\ell(I)$: han de
cortarse.)}
\end{exercise}

\begin{exercise}[$\star\star\star$]\label{exo:b3:measure:9}
Demostrar que todo $A \in \mathcal L$ con $\lambda(A) > 0$ contiene un
subconjunto no medible. \emph{(Interséquese $A$ con los trasladados
$V + q$ del conjunto de Vitali: si todos los $A \cap (V+q)$ fueran
medibles, cada uno sería nulo por el argumento
del~\cref{thm:b3:measure:vitali} --- Steinhaus
(el~\cref{exo:b3:measure:8}) ayuda: un conjunto medible de
\omterm{def:b3:measure:measure}{medida} positiva dentro de $V + q$ daría
a $(V+q) - (V+q) \supseteq$ un intervalo, en contra de que ese conjunto
de diferencias corta a $\Q$ solo en $0$; concluir con la
subaditividad.)}
\end{exercise}

\begin{exercise}[$\star\star$]\label{exo:b3:measure:10}
Demostrar que $A \subseteq \R$ con $\lambda^*(A) < \infty$ es medible
Lebesgue si y solo si para todo $\varepsilon > 0$ existe un abierto
$U \supseteq A$ con $\lambda^*(U \setminus A) <
\varepsilon$, si y solo si existe un conjunto $G_\delta$ $G \supseteq A$
con $\lambda^*(G\setminus A) = 0$. (Así, los conjuntos de Lebesgue son
los de Borel módulo conjuntos nulos.)
\end{exercise}

\begin{exercise}[$\star\star$]\label{exo:b3:measure:11}
(\omterm{def:b3:topology:continuity}{Continuidad} a lo largo de límites
monótonos, y su optimalidad) (a) Demostrar que, para conjuntos medibles,
$\mu(\liminf A_n) \leq
\liminf\mu(A_n)$ (Fatou para conjuntos) y que, si
$\mu\bigl(\bigcup A_n\bigr) < \infty$, también
$\limsup\mu(A_n) \leq \mu(\limsup A_n)$.
(b) Exhibir, para la \omterm{def:b3:measure:lebesgueouter}{medida de
Lebesgue} sobre $\R$, una sucesión con $\mu(A_n) = 1$ para todo $n$ y,
sin embargo, $\mu(\limsup A_n) = 0$: la hipótesis de finitud de la
segunda desigualdad no es decorativa. (c) Deducir: si
$\sum\mu(A_n) < \infty$, entonces $\mu(\limsup A_n) = 0$
(Borel--Cantelli de nuevo), y si los $A_n$ crecen o decrecen (con
$\mu(A_1) < \infty$ en el caso decreciente),
$\mu(\lim A_n) = \lim\mu(A_n)$.
\end{exercise}

\begin{exercise}[$\star\star\star$]\label{exo:b3:measure:12}
(Teorema de Egórov) Sean $\mu(X) < \infty$ y $f_n \to f$ puntualmente,
todas medibles (con valores reales). Para $k, N \geq 1$, póngase
\[
E_{k,N} = \bigcap_{n \geq N}\Bigl\{x : \abs{f_n(x) - f(x)}
\leq \tfrac1k\Bigr\} .
\]
(a) Demostrar que, para $k$ fijo, $E_{k,N} \nearrow X$ cuando $N \to
\infty$, y deducir $N_k$ con $\mu(X \setminus E_{k,N_k})
\leq \varepsilon2^{-k}$. (b) Concluir el \emph{teorema de Egórov}: para
todo $\varepsilon
> 0$ existe un medible $A$ con $\mu(X\setminus A) \leq
\varepsilon$ tal que $f_n \to f$ \emph{uniformemente sobre $A$} --- la
convergencia puntual es convergencia uniforme salvo en un conjunto
arbitrariamente pequeño. (c) Demostrar que el teorema falla en
$(\R, \lambda)$: las jorobas móviles $f_n = \mathbf 1_{\intcc n{n+1}}$
convergen puntualmente a $0$ pero uniformemente en el complementario de
ningún conjunto de \omterm{def:b3:measure:measure}{medida} finita.
¿Dónde usó (a) que $\mu(X) < \infty$?
\end{exercise}

\section{Problema: la escalera de Cantor--Vitali y un
conjunto medible que no es de Borel}

\begin{omfigure}
\begin{tikzpicture}[scale=5.2]
  \draw[->] (0,0) -- (1.06,0) node[right, font=\small] {$x$};
  \draw[->] (0,0) -- (0,1.06);
  \draw[black!30] (1,0) -- (1,1) -- (0,1);
  \draw[omThm, thick]
    (0,0)
    -- (0.037,0.125) -- (0.074,0.125)
    -- (0.111,0.25)  -- (0.222,0.25)
    -- (0.259,0.375) -- (0.296,0.375)
    -- (0.333,0.5)   -- (0.667,0.5)
    -- (0.704,0.625) -- (0.741,0.625)
    -- (0.778,0.75)  -- (0.889,0.75)
    -- (0.926,0.875) -- (0.963,0.875)
    -- (1,1);
  \node[font=\small, omThm] at (0.32,0.86)
    {$c$ (aproximación de profundidad 3)};
\end{tikzpicture}

{\small La escalera de Cantor--Vitali: constante en cada hueco del
conjunto de Cantor y, aun así, subiendo de $0$ a $1$ de manera
\omterm{def:b3:topology:continuity}{continua}. Su derivada se anula en
casi todo punto --- toda la subida ocurre sobre un conjunto nulo.}
\end{omfigure}

\begin{problem}[{Problema de fin de semana --- la escalera del diablo
y $\mathcal B(\R) \subsetneq \mathcal L$}]
\label{pb:b3:measure:1}
Construiremos la \emph{función de Cantor--Vitali} (la escalera del
diablo), la usaremos para transportar la
\omterm{def:b3:measure:measure}{medida} de manera patológica y
concluiremos con un teorema que ningún argumento blando proporciona:
existen conjuntos medibles Lebesgue que no son de Borel. Notación: $C$
es el conjunto de Cantor, $C_n$ su etapa $n$-ésima ($2^n$ intervalos de
longitud $3^{-n}$), y todo $x \in C$ tiene dígitos ternarios
$x = \sum 2b_n3^{-n}$, $b_n \in \{0,1\}$
(\cref{exo:b3:topology:10}).

\medskip
\noindent\textbf{Parte I --- La escalera.} Defínase $c_0(x) = x$ y
$c_{n+1}$ a partir de $c_n$ mediante
\[
c_{n+1}(x) = \begin{cases}
\tfrac12\,c_n(3x) & 0 \leq x \leq \tfrac13,\\[2pt]
\tfrac12 & \tfrac13 \leq x \leq \tfrac23,\\[2pt]
\tfrac12 + \tfrac12\,c_n(3x - 2) & \tfrac23 \leq x \leq 1.
\end{cases}
\]
\begin{enumerate}
  \item Demostrar que cada $c_n$ es
  \omterm{def:b3:topology:continuity}{continua} y no decreciente, con
  $c_n(0) = 0$, $c_n(1) = 1$, y que
        $\norm{c_{n+1} - c_n}_\infty \leq
        \tfrac12\norm{c_n - c_{n-1}}_\infty$.
  \item Deducir que $(c_n)$ converge uniformemente a una $c$
  \omterm{def:b3:topology:continuity}{continua} y no decreciente con
  $c(0) = 0$, $c(1) = 1$ (\emph{la} función de Cantor--Vitali), que
  cumple las mismas relaciones autosemejantes que los $c_{n+1}$
  anteriores. \item Demostrar que $c$ es constante en cada
  \omterm{def:b3:topology:components}{componente conexa} de
  $\intcc01\setminus C$ y que, para $x = \sum_n
        2b_n3^{-n} \in C$, $c(x) = \sum_n b_n2^{-n}$: la escalera lee en
  binario los dígitos de Cantor (la función $g$
  del~\cref{pb:b3:topology:1}, hecha monótona y global). \item Deducir
  que $c$ es derivable, con $c' = 0$, en todo punto de
  $\intcc01\setminus C$: $c' = 0$ \emph{$\lambda$-en casi todo punto}
  (el~\cref{ex:b3:measure:cantor}). Concluir que el teorema fundamental
  del cálculo falla para $c$:
        \[
        c(1) - c(0) = 1 \neq 0 = \int_0^1 c'(t)\,\dd t
        \]
        (donde la integral se toma sobre el conjunto de
        \omterm{def:b3:measure:measure}{medida} total en el que
        $c' = 0$; anticipando el~\cref{ch:b3:lebesgue}, los conjuntos
        nulos no afectan a las integrales). ¿Qué hipótesis del teorema
        fundamental de $\mathcal C^1$ se viola? \item Demostrar que
        $c(C) = \intcc01$: el conjunto nulo $C$ se aplica \emph{sobre}
        un conjunto de \omterm{def:b3:measure:measure}{medida} total.
\end{enumerate}

\medskip
\noindent\textbf{Parte II --- El
\omterm{def:b3:topology:continuity}{homeomorfismo} torcido.} Sea
$h(x) = \frac{c(x) + x}{2}$.
\begin{enumerate}[resume]
  \item Demostrar que $h \colon \intcc01 \to \intcc01$ es un
  \omterm{def:b3:topology:continuity}{homeomorfismo} (estrictamente
  creciente, \omterm{def:b3:topology:continuity}{continua} y
  sobreyectiva). \item Demostrar que
  $\lambda\bigl(h(\intcc01\setminus C)\bigr) =
        \tfrac12$: sobre cada hueco de longitud $\ell$, $h$ actúa como
  una aplicación afín de pendiente $\tfrac12$, y los huecos tienen
  longitud total $1$. \item Deducir
  $\lambda\bigl(h(C)\bigr) = \tfrac12$: la imagen
  \omterm{def:b3:topology:continuity}{homeomorfa} de un conjunto nulo
  puede tener \omterm{def:b3:measure:measure}{medida} positiva. (¿En qué
  contradice esto la intuición ingenua sobre el «tamaño»?)
\end{enumerate}

\medskip
\noindent\textbf{Parte III --- Un conjunto medible que no es de Borel.}
\begin{enumerate}[resume]
  \item Por el~\cref{exo:b3:measure:9}, tómese un $W
        \subseteq h(C)$ no medible. Demostrar que
  $Z = h^{-1}(W) \subseteq C$ es medible Lebesgue. \emph{(Es un
  subconjunto de un conjunto nulo;
  \omterm{def:b3:complete:complete}{completitud},
  el~\cref{thm:b3:measure:caratheodory}.)} \item Demostrar que la
  preimagen de un conjunto de Borel por una
  \omterm{def:b3:topology:continuity}{aplicación continua} es de Borel.
  \emph{(Principio de los buenos conjuntos:
  $\{B : h^{-1}(B) \in \mathcal B\}$ es una $\sigma$-álgebra que
  contiene los abiertos --- atención
  al sentido de la aplicación.)} \item Concluir que $Z$ \emph{no} es de
  Borel: si lo fuera, $W = (h^{-1})^{-1}(Z)$ sería de Borel (aplíquese
  la pregunta 10 a la \omterm{def:b3:topology:continuity}{continua}
  $h^{-1}$) y, por tanto, medible --- contradicción. Por consiguiente,
        \[
        \boxed{\ \mathcal B(\R) \subsetneq \mathcal L\ }
        \]
        y la \omterm{def:b3:complete:complete}{completitud} amplía
        genuinamente el mundo de Borel. \item Exhibir una función $g$
        medible Lebesgue y una función $\varphi$
        \omterm{def:b3:topology:continuity}{continua} tales que $g \circ
        \varphi$ no sea medible Lebesgue: la medibilidad, a diferencia
        de la \omterm{def:b3:topology:continuity}{continuidad}, no se
        compone. \emph{(Tómense $g =
        \mathbf 1_Z$ y $\varphi = h^{-1}$, anticipando la definición de
        función medible del~\cref{ch:b3:lebesgue}: las preimágenes de
        conjuntos de Borel son conjuntos de Lebesgue. ¿Dónde hay que
        tener cuidado con qué $\sigma$-álgebra se usa en el destino?)}
\end{enumerate}

\medskip
\noindent\textbf{Parte IV --- Epílogo.}
\begin{enumerate}[resume]
  \item Ordenar por inclusión estricta las clases siguientes,
  justificando cada estrictud con un ejemplo de este capítulo o de su
  problema: los conjuntos numerables; los conjuntos de Borel nulos; los
  conjuntos nulos de Lebesgue; los conjuntos de Borel; los conjuntos de
  Lebesgue; los conjuntos arbitrarios.
\end{enumerate}

\medskip
\noindent\textbf{Parte V --- La \omterm{def:b3:measure:measure}{medida}
de Cantor: masa sobre un conjunto nulo.} La escalera es la función de
distribución de una \omterm{def:b3:measure:measure}{medida} notable, que
ahora construiremos con las herramientas propias de este capítulo.
\begin{enumerate}[resume]
  \item (Lebesgue--Stieltjes, existencia) Sea $F\colon\R\to\R$ no
  decreciente, \omterm{def:b3:topology:continuity}{continua} y acotada.
  Sobre los intervalos semiabiertos defínase
  $\rho\bigl(\intoc ab\bigr) = F(b) -
        F(a)$ y, para $A \subseteq \R$,
        \[
        \mu_F^*(A) = \inf\Bigl\{\sum_k\bigl(F(b_k) -
        F(a_k)\bigr) : A \subseteq
        \bigcup_k\intoc{a_k}{b_k}\Bigr\} .
        \]
        Demostrar que $\mu_F^*$ es una
        \omterm{def:b3:measure:outer}{medida exterior} y que
        $\mu_F^*\bigl(\intoc ab\bigr) = F(b) - F(a)$
        \emph{(imítese el argumento de \omterm{def:b3:topology:compact}{compacidad}
        del~\cref{thm:b3:measure:lebesgue}, agrandando cada
        $\intoc{a_k}{b_k}$ hasta un intervalo abierto con un coste $F$
        de $\leq \varepsilon2^{-k}$ --- ¿dónde se usa la
        \omterm{def:b3:topology:continuity}{continuidad} de $F$?)}.
        \item Demostrar que todo conjunto de Borel es $\mu_F^*$-medible
        en el sentido de Carathéodory \emph{(como en el caso de
        Lebesgue, basta comprobarlo sobre las semirrectas; sígase la
        demostración de la aplicación
        del~\cref{thm:b3:measure:caratheodory})}, de modo que
        $\mu_F = \mu_F^*$ restringida a $\mathcal B(\R)$ es una
        \omterm{def:b3:measure:measure}{medida} con
        $\mu_F(\intoc ab) = F(b) - F(a)$: la
        \emph{\omterm{def:b3:measure:measure}{medida} de
        Lebesgue--Stieltjes} de $F$. \item Aplíquese esto a la escalera
        ($F = c$ extendida por $0$ sobre $\R_-$ y por $1$ sobre
        $\intco1\infty$): la
        \emph{\omterm{def:b3:measure:measure}{medida} de Cantor} $\mu$.
        Demostrar que $\mu(\R) = 1$, que todo hueco del conjunto de
        Cantor es $\mu$-nulo ($c$ es allí constante) y concluir que
        \[
        \mu(C) = 1, \qquad \lambda(C) = 0 :
        \]
        $\mu$ y $\lambda$ viven sobre soportes disjuntos ($C$ y su
        complementario). Dos \omterm{def:b3:measure:measure}{medidas} en
        esta situación se llaman \emph{mutuamente singulares}, y se
        escribe $\mu \perp
        \lambda$. \item Demostrar que $\mu$ no tiene átomos:
        $\mu(\{x\}) = 0$ para todo $x$
        \emph{(\omterm{def:b3:topology:continuity}{continuidad} de
        $c$)}. Una \omterm{def:b3:measure:measure}{medida} de
        probabilidad sin átomos soportada por un
        \omterm{def:b3:topology:compact}{compacto} de \omterm{def:b3:measure:lebesgueouter}{medida de Lebesgue}
        nula: compárese con las únicas
        \omterm{def:b3:measure:measure}{medidas} vistas hasta ahora.
        \item (Lanzamiento de monedas disfrazado) Para una palabra
        $(\varepsilon_1,
        \dots, \varepsilon_m) \in \{0,1\}^m$, sea $C_{\varepsilon}$ el
        conjunto de los $x \in C$ cuyos dígitos ternarios cumplen
        $b_i(x) = \varepsilon_i$ para $i \leq m$ (una de las $2^m$
        piezas de Cantor de profundidad $m$). Demostrar que
        $\mu(C_\varepsilon) = 2^{-m}$ \emph{(la escalera sube $2^{-m}$ a
        lo largo de esa pieza: úsese la Parte I, pregunta 3)}. La
        \omterm{def:b3:measure:measure}{medida} de Cantor es la ley de
        una sucesión infinita de lanzamientos de moneda equilibrada
        leída en base tres --- el~\cref{ch:b3:probability} hará esto
        exacto. \item Demostrar la autosemejanza: para todo $A$ de
        Borel,
        \[
        \mu(A) = \tfrac12\,\mu(3A) +
        \tfrac12\,\mu(3A - 2) ,
        \]
        donde $3A - 2 = \{3x - 2 : x \in A\}$ \emph{(compruébese sobre
        los intervalos generadores $\intoc ab$ mediante las relaciones
        autosemejantes de $c$, e invóquese después la unicidad,
        \cref{thm:b3:measure:uniqueness})}.
  \item Demostrar que la reflexión $s(x) = 1 - x$ conserva $\mu$:
  $\mu(s(A)) = \mu(A)$ \emph{(mediante $c(1 - x) = 1 -
        c(x)$, que se sigue de la simetría de la construcción ---
  demuéstrese)}. \item Calcular los dos primeros momentos de $\mu$, es
  decir, de un punto aleatorio $X$ con ley $\mu$ (las integrales pueden
  tratarse como límites de sumas sobre las piezas de profundidad $m$,
  anticipando el~\cref{ch:b3:lebesgue}): la simetría da
  $\int x\,\dd\mu = \frac12$, y la autosemejanza da
        \[
        \int x^2\,\dd\mu = \frac38,
        \qquad\text{luego}\qquad
        \operatorname{Var}(X) = \frac18 .
        \]
        Compárese con la ley uniforme sobre $\intcc01$ (varianza
        $\frac1{12}$): la masa de Cantor, empujada hacia los bordes, se
        dispersa \emph{más}. \item Demostrar que el \emph{soporte}
        (\omterm{def:b3:topology:topology}{topológico}) de $\mu$ --- el
        menor cerrado de \omterm{def:b3:measure:measure}{medida} total
        --- es exactamente $C$. \item (Síntesis) La escalera $c$ es
        \omterm{def:b3:topology:continuity}{continua} y no decreciente
        y, sin embargo, incumple el teorema fundamental del cálculo
        (Parte I); la \omterm{def:b3:measure:measure}{medida} $\mu_c$ es
        una probabilidad, sin átomos, singular respecto de $\lambda$.
        Explíquese en un párrafo breve cómo estas son dos caras de un
        mismo fenómeno y enúnciese la moraleja general: las funciones no
        decrecientes corresponden a
        \omterm{def:b3:measure:measure}{medidas}
        ($F \leftrightarrow \mu_F$), la derivabilidad en casi todo punto
        corresponde a la «parte absolutamente
        \omterm{def:b3:topology:continuity}{continua}», y $c$ es el
        testigo estándar de que una $F$
        \omterm{def:b3:topology:continuity}{continua} puede no tener
        \emph{ninguna} parte absolutamente
        \omterm{def:b3:topology:continuity}{continua}. \item (El módulo
        de \omterm{def:b3:topology:continuity}{continuidad} óptimo) Sea
        $s =
        \frac{\ln 2}{\ln 3}$. Demostrar que $c$ es
        \omterm{def:b3:topology:continuity}{continua} de tipo Hölder con
        exponente $s$:
        \[
        \abs{c(x) - c(y)} \leq 4\,\abs{x - y}^{s}
        \qquad (x, y \in \intcc01),
        \]
        y que ningún exponente $t > s$ puede servir, ni siquiera
        localmente. Deducir la forma en términos de \omterm{def:b3:measure:measure}{medida}: para todos
        $x$ y $r \in \intoc01$,
        \[
        \mu\bigl(\intcc{x - r}{x + r}\bigr) \leq 8\,r^{s} .
        \]
        \emph{(Compárese una malla triádica de profundidad $m$ con la
        escala de $\abs{x - y}$; la pregunta 18 da la subida a lo largo
        de cada pieza. El exponente $s$ es la dimensión \omterm{def:b3:topology:hausdorff}{de
        Hausdorff} de $C$, como dirán
        cursos posteriores.)} \item (La autosemejanza caracteriza $\mu$)
        Demostrar el recíproco de la pregunta 19: si $\nu$ es una
        \omterm{def:b3:measure:measure}{medida} de probabilidad sobre
        $\mathcal B(\R)$ soportada por $\intcc01$ y cumple
        \[
        \nu(A) = \tfrac12\,\nu(3A) + \tfrac12\,\nu(3A - 2)
        \qquad (A \in \mathcal B(\R)),
        \]
        entonces $\nu = \mu$. \emph{(Itérese la relación $m$ veces para
        repartir $\nu$ entre las $2^m$ piezas de Cantor de profundidad
        $m$, estímese $\nu(\intoc ab)$ contra el número de piezas
        contenidas en $\intoc ab$ y hágase $m \to \infty$; termínese con
        el~\cref{thm:b3:measure:uniqueness}.)}
\end{enumerate}
\end{problem}
